\documentclass[conference]{IEEEtran}
\IEEEoverridecommandlockouts

\usepackage[T1]{fontenc}
\usepackage[utf8]{inputenc}
\usepackage[brazil]{babel}
\usepackage{cite}
\usepackage{amsmath,amssymb,amsfonts}
\usepackage{graphicx}
\usepackage{textcomp}
\usepackage{xcolor}
\usepackage{url}
\usepackage{booktabs}
\usepackage{tabularx}
\usepackage{tikz}
\usetikzlibrary{arrows.meta, fit, backgrounds, positioning}

\def\BibTeX{{\rm B\kern-.05em{\sc i\kern-.025em b}\kern-.08em
    T\kern-.1667em\lower.7ex\hbox{E}\kern-.125emX}}

\begin{document}

\title{RPL como Protocolo de Roteamento para Backhaul Sem Fio em
Field Area Networks de Infraestruturas Críticas}

\author{
  \IEEEauthorblockN{Fernando Sabino Dantas}
  \IEEEauthorblockA{
    \textit{Programa de Pós-Graduação em Informática (PPGIa)}\\
    \textit{Pontifícia Universidade Católica do Paraná (PUCPR)}\\
    Curitiba, Brasil\\
    fsd.dantas@gmail.com
  }
}

\maketitle

% ─── ABSTRACT (English) ──────────────────────────────────────────────────────
\begin{abstract}
Modern critical infrastructures --- particularly electric distribution networks
--- require pervasive, resilient, and cost-effective communication across
geographically dispersed field assets. Traditional backhaul solutions (fiber,
cellular, power-line communication) present economic, operational, or technical
constraints that limit their applicability in last-mile field environments.
Wireless mesh Field Area Networks (FANs), operating over constrained radio
links, have emerged as a viable backhaul alternative, with RPL (IPv6 Routing
Protocol for Low-Power and Lossy Networks, RFC~6550) serving as their standard
routing layer. This paper presents an overview and tutorial of RPL in the
context of FAN backhaul, covering its DODAG-based architecture, control
messaging (DIO, DAO, DIS), objective functions, routing metrics, and resilience
mechanisms for link failure and dynamic topology reconfiguration. The Wi-SUN
FAN profile, which formalizes RPL deployment for utility networks, is discussed
alongside simulation tools (Contiki-NG, Cooja) and applications in advanced
metering infrastructure, distribution automation, and wireless SCADA. Open
challenges in scalability, determinism, and integration with operational
technology protocols (IEC~61850, DNP3) are also addressed.
\end{abstract}

\begin{IEEEkeywords}
RPL, FAN, backhaul sem fio, LLN, DODAG, smart grid, automação de distribuição,
Wi-SUN, infraestrutura crítica
\end{IEEEkeywords}

% ─── I. INTRODUÇÃO ───────────────────────────────────────────────────────────
\section{Introdução}

Redes de distribuição elétrica, redes de água e gás, e infraestruturas de
transporte urbano dependem de comunicação confiável entre ativos de campo
dispersos geograficamente --- medidores, sensores, chaves automatizadas,
transformadores e subestações. A ausência de comunicação ou sua degradação
compromete diretamente a operação do sistema: sem telemetria em tempo real,
faltas não são isoladas automaticamente, o balanço de carga é realizado de
forma manual e a qualidade de energia não pode ser monitorada continuamente
\cite{gungor2011smartgrid}.

\subsection{Limitações das Soluções de Backhaul Convencionais}

Três tecnologias dominam o backhaul em infraestruturas críticas hoje:

\textbf{Fibra óptica}: oferece largura de banda e latência ideais, mas seu custo
de implantação em campo (obras civis, passagem por estruturas de terceiros) é
proibitivo para cobrir a capilaridade total de uma rede de distribuição. Em
áreas rurais, a ausência de dutos existentes eleva o custo por km a valores
incompatíveis com o orçamento de operação e manutenção~\cite{gungor2011smartgrid}.

\textbf{LTE privado (Band 31 / 450~MHz)}: concessionárias brasileiras operam
redes LTE próprias na faixa de 450~MHz, licenciada pela ANATEL para serviços
de infraestrutura crítica. Essa faixa sub-GHz oferece cobertura de área
superior às bandas convencionais e elimina a dependência operacional de
terceiros~\cite{saputro2012routing}. Contudo, o LTE privado atua como WAN
entre concentradores e o sistema de gestão; a camada de campo --- milhares de
medidores, sensores e atuadores com custo unitário baixo e consumo restrito
--- demanda uma rede de acesso com custo de módulo e energia inferiores ao de
uma interface LTE dedicada por dispositivo. A malha RPL/FAN opera exatamente
nessa camada, agregando o tráfego de campo nos concentradores que então se
comunicam com o head-end via LTE Band~31. As duas tecnologias são
complementares, não concorrentes.

\textbf{Comunicação por linha de potência (PLC)}: aproveita a infraestrutura
elétrica existente, mas sofre com interferência de alta frequência gerada pelos
próprios equipamentos da rede, atenuação variável com a carga e dificuldade de
operação em topologias de média tensão.

\subsection{Field Area Networks como Alternativa de Backhaul}

Redes de campo sem fio --- denominadas Field Area Networks (FANs) ou
Neighborhood Area Networks (NANs) no contexto de medição avançada --- operam na
camada entre os dispositivos de campo (medidores, sensores, atuadores) e os
concentradores de dados ou gateways que se conectam à rede de gestão. A
Figura~\ref{fig:fan} ilustra a hierarquia típica: dispositivos de campo formam
uma malha sem fio multihop que converge em concentradores de dados, os quais se
conectam ao sistema de back-office via WAN~\cite{ancillotti2013rpl}.

\begin{figure}[t]
  \centering
  \resizebox{\columnwidth}{!}{%
  \begin{tikzpicture}[
    dev/.style   = {circle, draw=black!55, fill=gray!25,
                    minimum size=9pt, inner sep=0},
    infra/.style = {rectangle, rounded corners=3pt,
                    draw=black!75, fill=white,
                    minimum width=2.2cm, minimum height=1.0cm,
                    font=\normalsize\bfseries, align=center},
    rpl/.style   = {-Latex, black!40, thin,
                    shorten <=2pt, shorten >=2pt},
    wan/.style   = {-Latex, black!80, thick, dashed,
                    shorten <=5pt, shorten >=5pt},
    >=Latex
  ]
    %% ── field devices ───────────────────────────────────────
    \node[dev] (d1) at (0.0,  0.0) {};
    \node[dev] (d2) at (0.0,  1.2) {};
    \node[dev] (d3) at (1.0,  1.9) {};
    \node[dev] (d4) at (1.0,  0.6) {};
    \node[dev] (d5) at (0.5, -0.7) {};
    \node[dev] (d6) at (2.0,  1.2) {};

    %% ── infrastructure ──────────────────────────────────────
    \node[infra] (dap) at (4.0, 0.6) {Concentrador\\(DAP)};
    \node[infra] (he)  at (7.4, 0.6) {Head-end\\SCADA};

    %% ── RPL upward routes ───────────────────────────────────
    \draw[rpl] (d1) -- (d4);
    \draw[rpl] (d2) -- (d4);
    \draw[rpl] (d5) -- (d4);
    \draw[rpl] (d3) -- (d6);
    \draw[rpl] (d6) -- (dap);
    \draw[rpl] (d4) -- (dap);

    %% ── WAN backhaul ────────────────────────────────────────
    \draw[wan] (dap) --
      node[above, font=\normalsize] {LTE Band~31 (450\,MHz)}
      (he);

    %% ── FAN zone background ─────────────────────────────────
    \begin{scope}[on background layer]
      \fill[gray!10, rounded corners=7pt]
        (-0.6,-1.1) rectangle (3.0, 2.6);
    \end{scope}

    %% ── zone labels ─────────────────────────────────────────
    \node[font=\normalsize, text=black!55, align=center]
      at (1.2, -1.7) {FAN --- RPL / IEEE~802.15.4g / Wi-SUN};
    \node[font=\normalsize, text=black!55]
      at (5.7, -1.7) {WAN};

    %% ── device legend ───────────────────────────────────────
    \node[dev, scale=0.85] at (-0.55, 3.1) {};
    \node[font=\normalsize, text=black!60, right]
      at (-0.3, 3.1) {Medidor / Sensor / Atuador};
  \end{tikzpicture}%
  }
  \caption{Arquitetura de comunicação em redes de distribuição.
           O RPL atua na camada FAN, provendo roteamento multihop
           entre dispositivos de campo e concentradores.}
  \label{fig:fan}
\end{figure}

A FAN opera tipicamente sobre IEEE~802.15.4g~\cite{ieee802154g} --- extensão
do padrão original para comunicação em redes utilitárias, com suporte a
frequências sub-GHz (169, 433, 868/915~MHz) que oferecem maior alcance (500~m
a 2~km em espaço aberto) e melhor penetração em estruturas físicas do que as
faixas de 2,4~GHz. Sobre essa camada física, o padrão Wi-SUN~\cite{wisun2022}
especifica a pilha completa: IEEE~802.15.4g / 6LoWPAN / IPv6 / RPL, com perfis
de segurança e interoperabilidade para implantação em escala por concessionárias.

\subsection{RPL como Protocolo de Roteamento para FANs}

O RPL (\textit{IPv6 Routing Protocol for Low-Power and Lossy
Networks})~\cite{winter2012rpl}, padronizado pelo IETF na RFC~6550, foi
concebido especificamente para redes com enlace sem fio de baixa confiabilidade
e dispositivos com recursos restritos --- as mesmas condições encontradas em
FANs de infraestrutura. Ao contrário de protocolos tradicionais (AODV, OSPF,
DSR), o RPL considera métricas de qualidade de enlace e energia na construção
de rotas, suporta topologias com milhares de nós e opera com overhead de
controle compatível com dispositivos embarcados~\cite{ancillotti2013rpl}.

Este trabalho apresenta o RPL no contexto de backhaul FAN, organizado da
seguinte forma: a Seção~\ref{sec:arch} descreve a arquitetura e o funcionamento
do protocolo; a Seção~\ref{sec:resilience} analisa os mecanismos de resiliência;
a Seção~\ref{sec:tools} discute ferramentas de simulação; a
Seção~\ref{sec:apps} aborda aplicações em infraestruturas críticas; e a
Seção~\ref{sec:conclusion} apresenta as conclusões.

% ─── II. ARQUITETURA E FUNCIONAMENTO ─────────────────────────────────────────
\section{Arquitetura e Funcionamento do RPL}
\label{sec:arch}

\subsection{Estrutura DODAG}

O RPL organiza a topologia em grafos acíclicos dirigidos orientados a destino
--- DODAGs (\textit{Destination Oriented Directed Acyclic Graphs}). Em uma FAN,
o \textit{DODAG Root} é tipicamente o concentrador de dados ou data aggregation
point (DAP) que conecta a malha de campo à WAN. Cada nó mantém um valor de
\textit{rank}, representando sua posição relativa ao concentrador: o concentrador
possui rank mínimo e dispositivos mais distantes possuem rank crescente.

O identificador de cada instância DODAG é o \textit{DODAGID}, acompanhado por
um \textit{DODAGVersionNumber} incrementado a cada reconstrução global da
topologia. Múltiplas instâncias RPL podem coexistir na mesma malha física,
permitindo, por exemplo, que o tráfego de medição e o tráfego de automação de
distribuição utilizem topologias otimizadas para métricas
distintas~\cite{winter2012rpl}.

\subsection{Mensagens de Controle}

O RPL define três tipos de mensagens ICMPv6~\cite{winter2012rpl}:

\textbf{DIO (\textit{DODAG Information Object})}: Propagado pelo concentrador e
retransmitido pelos nós descendentes, o DIO contém DODAGID, versão, rank do
transmissor, função objetivo em uso e parâmetros de configuração. Dispositivos
de campo utilizam o DIO para selecionar o pai preferencial e calcular seu rank.

\textbf{DAO (\textit{Destination Advertisement Object})}: Propaga informações de
alcançabilidade dos dispositivos de campo em direção ao concentrador, habilitando
o roteamento descendente (do concentrador para os dispositivos). Em modo
\textit{storing}, cada nó retém tabelas de roteamento dos seus descendentes; em
modo \textit{non-storing}, somente o concentrador mantém o mapa completo e utiliza
roteamento por fonte para tráfego descendente.

\textbf{DIS (\textit{DODAG Information Solicitation})}: Enviado por dispositivos
recém-implantados ou que perderam conectividade para acelerar a (re)descoberta da
topologia. Ao receber um DIS, os vizinhos respondem com DIO, reduzindo o tempo de
ingresso na rede --- relevante em cenários de substituição de equipamento em campo.

\subsection{Algoritmo Trickle}

O RPL utiliza o algoritmo Trickle~\cite{rfc6206} para controlar a frequência de
transmissão dos DIOs. O Trickle adapta o intervalo entre limites mínimo
($I_{min}$) e máximo ($I_{max}$): em topologias estáveis, o intervalo dobra a
cada ciclo, reduzindo drasticamente o overhead de controle; quando inconsistências
são detectadas, o intervalo reinicia a $I_{min}$ para convergência rápida. Em uma
FAN com centenas de medidores, este mecanismo é determinante para manter o canal
de rádio disponível para tráfego de dados.

\subsection{Funções Objetivo e Métricas}

A Função Objetivo (OF) define como o rank é calculado e como o pai preferencial
é selecionado. As duas funções padronizadas são:

\textbf{OF0}~\cite{rfc6552}: Utiliza hop count como métrica, minimizando a
profundidade da árvore. Adequada para topologias homogêneas onde a qualidade
de enlace é relativamente uniforme.

\textbf{MRHOF}~\cite{rfc6719}: Utiliza ETX (\textit{Expected Transmission
Count}) como métrica primária, com histerese para evitar oscilação entre pais
de qualidade similar. O ETX captura tanto a taxa de perda de pacotes quanto as
retransmissões de camada de enlace --- métrica especialmente relevante em
ambientes com interferência eletromagnética variável, como subestações e
corredores de linhas de transmissão.

O framework de métricas~\cite{rfc6551} suporta adicionalmente energia residual,
latência e confiabilidade de rota. A Tabela~\ref{tab:metrics} resume as opções
disponíveis e sua relevância para FANs.

\begin{table}[t]
  \caption{Métricas de Roteamento RPL e Relevância para FAN}
  \label{tab:metrics}
  \centering
  \begin{tabularx}{\columnwidth}{lX}
    \toprule
    \textbf{Métrica} & \textbf{Relevância em FAN} \\
    \midrule
    Hop count    & Minimiza latência em topologias homogêneas \\
    ETX          & Captura qualidade real do enlace em campo \\
    Energia      & Estende vida útil de dispositivos a bateria \\
    Latência     & Requisito de DA e proteção de distância \\
    Confiabilidade & Garante entrega em canais com interferência \\
    \bottomrule
  \end{tabularx}
\end{table}

\subsection{Modos de Operação}

\textbf{Modo Storing}: Cada nó mantém tabela de roteamento para todos os
descendentes; encaminhamento é hop-by-hop. Adequado para FANs de médio porte
com concentradores de capacidade razoável.

\textbf{Modo Non-Storing}: Somente o concentrador mantém informações de
roteamento; tráfego descendente usa roteamento por fonte (path inserido no
cabeçalho IPv6). Adequado para dispositivos de campo com memória mínima, ao
custo de maior tamanho dos pacotes e dependência do concentrador no
encaminhamento~\cite{winter2012rpl}.

% ─── III. RESILIÊNCIA ────────────────────────────────────────────────────────
\section{Resiliência no RPL}
\label{sec:resilience}

\subsection{Detecção de Falhas de Enlace}

Em ambientes de campo, enlaces sem fio estão sujeitos a degradação por
interferência eletromagnética, obstrução física (vegetação, estruturas metálicas)
e variações atmosféricas. O RPL detecta falhas em múltiplas camadas: na camada
de enlace, o IEEE~802.15.4 sinaliza falha via ausência de ACK após o número
máximo de retransmissões; na camada de rede, a ausência de DIOs do pai
preferencial por múltiplos intervalos Trickle indica falha do nó pai.

Dispositivos monitoram continuamente o ETX dos enlaces com candidatos a pai.
Degradação abrupta do ETX --- típica de equipamentos com interferência
variável --- dispara reavaliação do pai sem aguardar detecção completa de
falha~\cite{accettura2011rpl}.

\subsection{Reparo Local e Reconfiguração Dinâmica}

O reparo local é acionado quando um dispositivo perde conectividade com seu pai
preferencial. O dispositivo examina sua lista de pais alternativos registrados
a partir de DIOs previamente recebidos. Se um pai alternativo com rank inferior
ao do dispositivo existir, a troca é realizada localmente, sem propagação à
raiz, minimizando perturbação na malha~\cite{winter2012rpl}.

Este mecanismo é crítico em FANs de distribuição elétrica: durante uma falta
de linha, múltiplos dispositivos de campo podem perder enlace simultaneamente.
A reconfiguração local paralela e independente de cada dispositivo garante
continuidade do fluxo de telemetria sem coordenação centralizada --- propriedade
especialmente relevante em cenários de black-out parcial onde o próprio
concentrador pode estar temporariamente inacessível.

A detecção de inconsistência de rank previne loops: se um nó recebe DIO de um
suposto descendente com rank inferior ao seu, detecta inversão de rota e
descarta o pai atual imediatamente~\cite{winter2012rpl}.

\subsection{Reparo Global}

Quando o reparo local é insuficiente --- particionamento da rede ou degradação
sistêmica após evento de grande escala --- o concentrador incrementa o
\textit{DODAGVersionNumber}, forçando reconstrução completa do DODAG. O
algoritmo Trickle reiniciado a $I_{min}$ propaga o novo DIO rapidamente pela
malha. O custo de convergência é elevado, devendo ser reservado para falhas que
o reparo local não pode corrigir.

\subsection{Segurança}

O RPL define modos de segurança com suporte a AES-128 CCM* para proteção de
mensagens de controle. Em FANs de infraestrutura crítica, ataques documentados
incluem: o \textit{rank attack} (nó malicioso anuncia rank baixo para atrair
tráfego), o \textit{version number attack} (força reparos globais contínuos,
degradando disponibilidade) e o \textit{sinkhole attack} (nó comprometido
intercepta tráfego da malha)~\cite{korte2012rplattacks}. O perfil Wi-SUN
endereça essas vulnerabilidades com autenticação mútua por certificados X.509 e
gerenciamento de chaves por KMP (\textit{Key Management Protocol}).

% ─── IV. FERRAMENTAS E FRAMEWORKS ────────────────────────────────────────────
\section{Ferramentas e Frameworks}
\label{sec:tools}

\subsection{Contiki-NG}

Contiki-NG~\cite{dunkels2004contiki} é o sistema operacional de referência para
desenvolvimento e avaliação de RPL. Sua arquitetura baseada em protothreads
permite concorrência cooperativa com consumo de 10--30~KB de flash e 2--5~KB de
RAM. Inclui implementações completas de RPL (storing e non-storing), 6LoWPAN,
TSCH e CoAP, cobrindo toda a pilha de um nó Wi-SUN. A compatibilidade com
plataformas como Zolertia RE-Mote (sub-GHz) e nRF52840 viabiliza prototipagem
direta para hardware de campo.

\subsection{Cooja}

O simulador Cooja~\cite{tsiftes2010cooja} emula redes com centenas de nós em
nível de instrução, executando o mesmo binário compilado para hardware real.
Para pesquisa em FANs, as funcionalidades mais relevantes são: visualização
gráfica da topologia DODAG ao longo do tempo, captura de pacotes integrada ao
Wireshark (IEEE~802.15.4/6LoWPAN/RPL), coleta de métricas de energia via
PowerTrace e scripting JavaScript para automação de experimentos com injeção
de falhas controlada.

\subsection{RIOT OS e ns-3}

O RIOT OS oferece modelo de threading POSIX com amplo suporte a hardware
embarcado e implementação do RPL conforme RFC~6550, adequada para avaliações de
conformidade. O ns-3 fornece modelos de canal de alta fidelidade para avaliação
em condições específicas de propagação, sendo utilizado em estudos de desempenho
em ambientes com interferência industrial ou obstáculos físicos densos --- cenários
típicos de subestações e corredores de linha de distribuição.

\subsection{Análise de Tráfego em Campo}

A validação de implantações reais demanda captura passiva do espectro sem
interferência na rede operacional. Hardware como o Zolertia ZOUL e adaptadores
ATUSB permitem análise de tráfego IEEE~802.15.4 via Wireshark, com dissectores
para 6LoWPAN e RPL. Em ambientes sub-GHz (Wi-SUN), analisadores de espectro com
suporte à faixa de 915~MHz são necessários para diagnóstico de coexistência com
outros equipamentos de campo.

% ─── V. APLICAÇÕES ───────────────────────────────────────────────────────────
\section{Aplicações em Infraestruturas Críticas}
\label{sec:apps}

\subsection{Infraestrutura de Medição Avançada (AMI)}

A AMI (\textit{Advanced Metering Infrastructure}) é a aplicação mais madura do
RPL em backhaul de campo. Medidores inteligentes, instalados em transformadores
de distribuição e ramais residenciais/comerciais, formam a malha FAN que
transporta leituras de consumo, alertas de qualidade de energia e comandos de
corte/religamento ao sistema de gestão da concessionária. O RPL, operando sobre
IEEE~802.15.4g~\cite{ieee802154g}, viabiliza a formação automática da topologia
multihop sem configuração manual por ponto~\cite{ancillotti2013rpl}.

A escala típica de uma AMI metropolitana --- dezenas de milhares de medidores
por concentrador --- exige as propriedades de escalabilidade e auto-organização
do RPL. O \textit{rank-based} forwarding elimina a necessidade de tabelas de
roteamento completas nos medidores, e o algoritmo Trickle garante que o overhead
de controle não degrada o canal à medida que a rede cresce.

\subsection{Automação de Distribuição e FDIR}

A automação de distribuição inclui chaveamento remoto, monitoramento de
transformadores e, especialmente, FDIR (\textit{Fault Detection, Isolation and
Restoration}) --- a capacidade de isolar automaticamente um segmento com falta
e restaurar o fornecimento pelos circuitos alternativos disponíveis. Para FDIR,
a latência de comunicação é crítica: o intervalo entre a detecção da falta e a
atuação da chave deve ser inferior a centenas de milissegundos em sistemas de
proteção de distância~\cite{gungor2011smartgrid}.

RPL puro não oferece garantias temporais determinísticas, mas sua integração com
IEEE~802.15.4e TSCH no stack 6TiSCH~\cite{de2014orchestra} acrescenta
escalonamento de slots que reduz jitter a valores compatíveis com automação de
subestação. Essa combinação posiciona RPL/6TiSCH como substituto sem fio para
enlaces SCADA cabeados em pontos onde a instalação física é inviável.

\subsection{Monitoramento de Qualidade de Energia e SCADA Sem Fio}

Sensores de qualidade de energia --- medidores de harmônicos, distorção de
tensão, fator de potência e afundamentos --- distribuídos ao longo de
alimentadores formam redes de monitoramento de baixo custo quando comunicam via
RPL/FAN. A telemetria contínua permite ao operador identificar clientes com
equipamentos geradores de harmônica, precificar energia reativa e antecipar
falhas de transformador~\cite{ancillotti2013rpl}.

O backhaul sem fio via RPL substitui com vantagem econômica os enlaces ponto-a-
ponto dedicados (radio UHF, fibra) em subestações secundárias e pontos de
monitoramento remotos, onde o volume de dados é baixo mas a disponibilidade
do enlace é mandatória.

\subsection{Desafios Abertos}

\textbf{Escalabilidade}: Em FANs com dezenas de milhares de nós, o custo de
reparo global após eventos de larga escala (tempestades que afetam centenas de
km de linha) pode sobrecarregar o canal de rádio. Pesquisas em RPL hierárquico
e particionamento de DODAG endereçam esse problema~\cite{accettura2011rpl}.

\textbf{Integração com OT}: Sistemas SCADA legados operam com protocolos como
DNP3 e IEC~61850, que requerem mapeamento sobre IPv6/UDP para transporte via
RPL. A coexistência de requisitos de latência determinística (proteção) e
coleta periódica (medição) na mesma malha FAN é um desafio de engenharia ativo.

\textbf{Interoperabilidade multi-vendor}: O perfil Wi-SUN~\cite{wisun2022}
avança na padronização, mas implantações mistas com equipamentos de diferentes
fabricantes ainda exigem validação extensiva de compatibilidade, especialmente
nos mecanismos de segurança e nos parâmetros Trickle.

% ─── VI. CONCLUSÃO ───────────────────────────────────────────────────────────
\section{Conclusão}
\label{sec:conclusion}

O RPL consolida-se como o protocolo de roteamento padrão para backhaul sem fio
em Field Area Networks de infraestruturas críticas. Sua arquitetura DODAG, aliada
ao algoritmo Trickle e às funções objetivo configuráveis, resolve os desafios de
roteamento em malhas de campo de forma energicamente eficiente e sem configuração
manual --- propriedades inatingíveis por protocolos projetados para redes
convencionais.

A resiliência do RPL diante de falhas de enlace e a reconfiguração dinâmica da
topologia são especialmente valiosas em redes de distribuição elétrica, onde
eventos físicos (faltas, intempéries) impactam simultaneamente múltiplos enlaces
de comunicação. O padrão Wi-SUN, ao formalizar o uso do RPL sobre IEEE~802.15.4g
com perfis de segurança e interoperabilidade, eleva o protocolo de proposta
acadêmica a componente operacional de AMIs em escala global.

As ferramentas Contiki-NG e Cooja fornecem ambiente completo para pesquisa e
prototipagem, viabilizando avaliação de desempenho antes da implantação em campo.
Os desafios remanescentes --- latência determinística para automação de
distribuição, escalabilidade de reparo em eventos de grande escala e integração
com protocolos OT legados --- representam linhas abertas de pesquisa com alto
potencial de impacto em infraestruturas de distribuição de próxima geração.

% ─── REFERÊNCIAS ─────────────────────────────────────────────────────────────
\bibliographystyle{IEEEtran}
\bibliography{references}

\end{document}
